\documentclass[11pt,letterpaper]{article}
\usepackage[margin=0.85in]{geometry}
\usepackage{lmodern}
\usepackage{hyperref}
\hypersetup{colorlinks=false, urlcolor=black, pdfborder=0 0 0}
\setlength{\parindent}{0pt}
\setlength{\parskip}{0.45em}
\raggedright
\pagestyle{empty}

\begin{document}

\noindent
{\bfseries Ke-Cheng Hsu (Clio)}\\
Ph.D.\ Candidate, Department of Economics, HKUST\\
(852) 6840\,9753 \quad $|$ \quad \href{mailto:khsu@connect.ust.hk}{khsu@connect.ust.hk}\\
\href{https://cliohsu1997.github.io/}{cliohsu1997.github.io} \quad $|$ \quad \href{https://github.com/cliohsu1997}{github.com/cliohsu1997}

\vspace{0.7em}
\noindent
October 21, 2026

\vspace{0.35em}
\noindent
Hiring Team\\
RBB Economics --- Australia\\
Re: \#26/27 Application --- Economist (Melbourne preferred)

\vspace{0.2em}
Dear Hiring Team,

I am writing to apply for the Economist position in RBB Economics' \#26/27 recruitment round, with Melbourne as my preferred Australian office (and I am also open to Sydney). I am a Ph.D.\ candidate in Economics at HKUST, expected to graduate in June 2027 (and able to finish earlier if helpful). My training centres on the three areas the posting asks for an interest in: microeconomics, industrial organisation, and econometrics. I am a native Mandarin speaker and professionally fluent in English. RBB appeals to me because competition economics is the whole of its practice, not a side line; I want to bring my empirical IO training to live matters on mergers, market investigations, abuse of dominance, vertical agreements, and pricing.

The posting describes four core tasks, and they map closely onto what I already do. First, researching and evaluating key industry dynamics, trends, and market structures: my job market paper studies how Taobao apparel sellers and buyers respond to return-policy changes across six categories, and I document category-level trends, return rates, and pricing patterns before turning to modelling. Second, programming, model building, and regression analysis: I clean large transaction data in Python, R, and C++, estimate a structural model of purchase and return behaviour by the generalized method of moments (GMM), and run regression and counterfactual analysis on HPC/SLURM and AWS. Third, applying quantitative and qualitative techniques: with Professors Kohei Kawaguchi (Advisor, UTokyo Economic Consulting Inc.) and Kei Kawai (Senior Consultant, Compass Lexecon), I combine descriptive evidence, regressions, and a dynamic oligopoly model to study pricing and inventory under resale price maintenance using large bookstore point-of-sale records. Fourth, assisting in the preparation of expert economic reports and presentations: I turn each project's results into figures and short written analyses on fixed deadlines, and I have lectured and TA'd for PhD and undergraduate economics (avg 4.72/5), so I am used to explaining technical material clearly.

On the requirements side, I hold an economics degree to Masters level (MPhil, HKUST) and am completing a Ph.D.; my Ph.D.-level IO coursework covered merger simulations, entry/exit counterfactuals, and firm-conduct identification. I write and present clearly, work well in teams, and learn new industry detail quickly, which I expect to matter when juggling multiple deadlines for multi-faceted and demanding clients.

I have attached my CV, this cover letter, and my transcripts. Thank you for your consideration.

\vspace{0.7em}
Sincerely,\\[1.0em]
Ke-Cheng Hsu (Clio)

\end{document}
